\begin{resumo}[Abstract]
\begin{otherlanguage*}{english}

Computational microbiome analysis requires integrating multiple
statistical tools, heterogeneous data formats, and complex
bioinformatic pipelines --- a scenario that hinders researchers
without a programming background. This work describes the design and
implementation of an academic web platform for integrated mycobiome
and metataxonomic analysis, developed to support three parallel
scientific projects: \textbf{INOVAHERB} (factorial ITS mycobiome
analysis via ANCOM-BC2), \textbf{Pós-Fogo} (16S post-fire soil
recovery time series via MaAsLin2), and \textbf{Biorremediação}
(microbiota correlation with edaphic variables). The architecture
adopts a microservices approach: a FastAPI (Python) backend exposing
a CQRS REST API, an asynchronous R Worker executing heavy statistical
analyses (DESeq2, ANCOM-BC2, MaAsLin2, SpiecEasi, Random Forest,
GSEA, FUNGuild, and PICRUSt2), a Next.js 14 frontend with interactive
visualisations via Plotly.js and Cytoscape.js, and a data layer
comprising PostgreSQL~16 (including Large Objects for binary storage) and Elasticsearch~8. Job scheduling
is handled through native PostgreSQL queues (\texttt{LISTEN/NOTIFY}
with \texttt{FOR UPDATE SKIP LOCKED}), without relying on external
message brokers. The platform is deployed on a lightweight Kubernetes
cluster (k3s) with five AMD nodes and supports institutional
authentication via Google OAuth restricted to the
\texttt{@ufvjm.edu.br} domain. Expected outcomes include a unified
dashboard with six PCoA plots, two per project, produced as the
central figure of the thesis upon completion of all three pipelines.

\vspace{\onelineskip}
\noindent
\textbf{Keywords}: bioinformatics; mycobiome; metataxonomics; web platform; FastAPI; R Worker; PostgreSQL; QIIME2.

\end{otherlanguage*}
\end{resumo}
